\section{Related Work}
\label{sec:related_work}

This section positions our thesis---that LLM checkpoint loading depends on workload shape and should be selected adaptively---against four bodies of work: serving systems, serverless cold-start systems, checkpoint formats and loading infrastructure, and operating-system I/O policy. The common gap is that prior work either optimises after weights are resident, avoids loads through caching or placement, or improves the kernel/storage substrate. Tensora measures the application-level load mechanics themselves.

\subsection{LLM Serving Systems}

Production LLM serving engines primarily optimise throughput and latency once weights are resident. vLLM introduces PagedAttention for efficient KV-cache management \cite{pagedattention}. Orca pioneers iteration-level batching \cite{orca}. Sarathi-Serve refines the prefill--decode scheduling trade-off \cite{sarathi_serve}. DeepSpeed Inference applies model parallelism and kernel fusion \cite{deepspeed_inference_paper}. These systems make the serving pipeline efficient after model initialisation; checkpoint loading is usually treated as a fixed prelude to the serving path.

Our vLLM experiments use a serving engine precisely to test whether backend effects survive that larger context. They do not replace serving-system research; they isolate a cost that serving systems usually inherit from the loader.

\subsection{Serverless Cold Start}

Serverless LLM serving makes cold start a first-order concern: idle models are evicted and must be reloaded on demand. ServerlessLLM \cite{FU:2024} proposes multi-tier loading with GPU/NVMe/DRAM placement and live migration. HydraServe \cite{hydraserve} minimises cold-start latency in public clouds through pipeline parallelism and placement optimisation. Snapshot-based warm-up approaches \cite{silva2021snapshot} pre-materialise process or model state to bypass part of startup.

These systems optimise \emph{whether}, \emph{where}, or \emph{when} loading occurs. Tensora asks a complementary question: when a load is unavoidable, which userspace submission policy moves the required bytes fastest for this checkpoint geometry? That question remains relevant even in tiered-storage systems because each tier still exposes an I/O workload to the loader.

\subsection{Checkpoint Formats and Loading Systems}

Checkpoint formats determine what byte ranges a loader must read and whether those reads are contiguous, sharded, partitioned, or fragmented. SafeTensors \cite{safetensors_docs} provides safe tensor serialisation with contiguous whole-file or multi-shard layouts. GGUF optimises for memory-mapped CPU inference and quantised deployments. ServerlessLLM's partitioned layout \cite{serverlessllm_repo} distributes tensors across partition files with a central index, enabling partial loading and tiered placement.

Recent systems work optimises the loading path from different layers. BlitzScale \cite{blitzscale} uses GPU-side model caching over NVLink and RDMA to share already-loaded models between inference instances. Most relevant to our work, PPC/MAIO \cite{fast26_maio} optimises model loading through a programmable kernel page cache, adapting prefetching and eviction to LLM I/O patterns. MAIO operates below the system-call interface; Tensora operates above it, comparing userspace I/O backends on identical logical read plans. These approaches are complementary: kernel cache policy determines how pages are served, while userspace backend policy determines how application reads are issued.

A separate line of characterization work studies LLM checkpoint/restore I/O patterns on parallel file systems \cite{arxiv2512_checkpoint_io}, covering buffered versus direct I/O paths and liburing performance. That work focuses on training checkpointing in HPC environments; our focus is inference checkpoint loading in single-node serving.

The gap we address is therefore format-conditioned backend selection: given a fixed checkpoint layout, which userspace backend and planning strategy should execute the read plan, and how does the answer change with model scale and integration layer?

\subsection{I/O Backend Policy}

The Linux I/O interface landscape spans blocking POSIX calls, asynchronous runtimes built above them, POSIX AIO, and \texttt{io\_uring}. Recent database research provides the closest methodological template for our work: Jasny et al. \cite{IOURING_DATABASE_GUIDELINES} show that \texttt{io\_uring} benefits depend on workload shape and queue depth. Bulk sequential reads may already amortise syscall overhead, while many-small or high-concurrency workloads benefit from batched submission.

LLM checkpoint loading spans both regimes. Contiguous SafeTensors shards resemble bulk sequential access with shard-level fanout. Partitioned ServerlessLLM layouts resemble grouped range workloads. Our work translates the workload-shape lesson to LLM loading and extends it with cross-layer evidence (Rust, Python, vLLM), capability-gated fallback when \texttt{io\_uring} is blocked, and an adaptive policy grounded in workload features rather than API names.

\subsection{Positioning}

We do not introduce a new checkpoint format, optimise steady-state serving throughput, or solve cluster scheduling. We measure load mechanics across four backend policies and two layouts, for six public checkpoints, at three integration layers. The contribution is operational: evidence that backend choice depends jointly on layout, scale, and availability, plus a format-aware adaptive policy with explicit overrides. This is the missing piece between format design (what bytes exist), serving engines (how loaded weights are used), and kernel-level storage optimisation (how pages are served).
